\chapter*{第三部分学习路线}
\addcontentsline{toc}{chapter}{第三部分学习路线：从电场、电路到电磁统一}
\markboth{第三部分学习路线}{第三部分学习路线}

\begin{learningbox}{从电场、电路到电磁统一}
第三部分进入电磁学。电场、磁场、电势、电流和磁通量都比较抽象，阅读时要坚持先看图、再看定义、最后看公式。
\end{learningbox}

\section*{一、这一部分的四层结构}

\begin{enumerate}[itemsep=0.2em]
  \item \textbf{电场层：} Day 15 到 Day 17 研究电荷在电场中的受力、做功和运动。
  \item \textbf{电路层：} Day 18 到 Day 20 研究电容、电流、电阻、电源和闭合电路。
  \item \textbf{磁场层：} Day 21 到 Day 22 研究磁场对电流和运动电荷的作用。
  \item \textbf{感应层：} Day 23 到 Day 24 研究变化的磁通量如何产生电。
\end{enumerate}

主线是：电场让电荷加速，运动电荷形成电流，电流产生磁场，变化的磁场又能产生电流。

\section*{二、阅读时抓住四个关键词}

\begin{itemize}[itemsep=0.2em]
  \item \textbf{方向：} 电场力、洛伦兹力、安培力、感应电流方向都要先画图。
  \item \textbf{能量：} 电场力做功、电源输出、电阻发热，本质上都在讨论能量转化。
  \item \textbf{等效：} 电容器、电源内阻、非纯电阻等模型常要化复杂为简单。
  \item \textbf{反馈：} 电磁感应中速度、电流、安培力常互相影响。
\end{itemize}

\section*{三、卡住时先回到哪里}

电场、电势题卡住，先回到“力做功与能量变化”；电路、磁场和感应题卡住，先画清电流、速度、磁场方向和磁通量变化。
